\documentclass[reqno]{amsart}
\usepackage[T1]{fontenc}
\usepackage{amsmath,amssymb,amsthm}
\usepackage{xr-hyper}
\usepackage[hidelinks]{hyperref}
\externaldocument[ec-]{20261124T223000--elliptic-curves-for-the-proof__workbook}
\externaldocument[rl-]{20261125T010000--rank-lemma-and-zero-estimate__workbook}
\externaldocument[dt-]{20261125T002000--definite-theta-lifts__workbook}
\externaldocument[bp-]{20261126T020210--bessel-poincare-series-and-their-gram-operators__workbook}

% Chapter number; \ref{chap} gives it.
\newcounter{chap}
% Chapter numbers, generated from 20261125T012000--grh-workbook__guide.org
% by scripts/make-chapters.py.  Do not edit: change the order of the
% index and rerun the script (./preview-site.sh and CI do this).
%
% Usage in a sheet, after \newcounter{chap}:
%   % Chapter numbers, generated from 20261125T012000--grh-workbook__guide.org
% by scripts/make-chapters.py.  Do not edit: change the order of the
% index and rerun the script (./preview-site.sh and CI do this).
%
% Usage in a sheet, after \newcounter{chap}:
%   % Chapter numbers, generated from 20261125T012000--grh-workbook__guide.org
% by scripts/make-chapters.py.  Do not edit: change the order of the
% index and rerun the script (./preview-site.sh and CI do this).
%
% Usage in a sheet, after \newcounter{chap}:
%   \input{chapters}\chapnumber{<sheet ID>}
\makeatletter
\providecommand\chapters@set[2]{\@namedef{chapters@#1}{#2}}
\chapters@set{20261125T001000}{1}% zeros-to-ramp-sums
\chapters@set{20261125T001100}{2}% inert-primes-and-estimate-b
\chapters@set{20261125T005000}{3}% galois-fourier-and-averaging
\chapters@set{20261125T005500}{4}% columns-from-a-sphere
\chapters@set{20261125T000000}{5}% slope-method
\chapters@set{20261126T030100}{6}% a-twist-for-each-prime
\chapters@set{20261125T003000}{7}% the-chain-and-the-square-root
\chapters@set{20261125T010000}{8}% rank-lemma-and-zero-estimate
\chapters@set{20261125T010500}{9}% poincare-series-and-near-orthogonality
\chapters@set{20261124T223000}{10}% elliptic-curves-for-the-proof
\chapters@set{20261125T009500}{11}% fock-model-and-theta-intertwiners
\chapters@set{20261124T223100}{12}% complex-multiplication-and-inert-reduction
\chapters@set{20261125T008000}{13}% one-cm-determinant
\chapters@set{20261124T214500}{14}% supersingular-theta-divisibility
\chapters@set{20261124T214600}{15}% the-isotropic-plane
\chapters@set{20261125T013700}{16}% pairings-on-a-formal-group-of-height-two
\chapters@set{20261125T013500}{17}% hecke-nondegeneracy-on-the-tree
\chapters@set{20261126T030200}{18}% the-bound-is-attained
\chapters@set{20261125T009000}{19}% kronecker-limit-and-chowla-selberg
\chapters@set{20261125T002000}{20}% definite-theta-lifts
\chapters@set{20261125T015500}{21}% siegel-weil-in-the-smallest-case
\chapters@set{20261125T002100}{22}% definite-theta-lifts-and-norms
\chapters@set{20261125T016500}{23}% hecke-operators-and-theta-lifts
\chapters@set{20261125T016700}{24}% divisibility-along-the-chain
\chapters@set{20261125T004000}{25}% denominators-and-sha
\chapters@set{20261125T004100}{26}% tate-twists-and-sha
\chapters@set{20261126T020430}{27}% the-minus-part-by-gauss-sums
\chapters@set{20261125T014100}{28}% congruence-modules
\chapters@set{20261125T014500}{29}% eisenstein-congruences-and-selmer-classes
\chapters@set{20261125T014700}{30}% pseudocharacters-and-the-length-bound
\chapters@set{20261125T011000}{31}% galois-representations-black-box
\chapters@set{20261126T030300}{32}% the-data
\chapters@set{20261126T020210}{33}% bessel-poincare-series-and-their-gram-operators
\chapters@set{20261126T030400}{34}% the-seesaw-identity
\chapters@set{20261126T030500}{35}% divisibility-of-the-composite-lift
\chapters@set{20261126T030600}{36}% theorem-a-short-forms
\chapters@set{20261126T030700}{37}% the-determinant-of-a-chain
\chapters@set{20261126T030800}{38}% the-cost-of-one-lift
\chapters@set{20261126T030900}{39}% denominators-of-projections-to-lifts
\chapters@set{20261126T031000}{40}% twists-with-few-congruences
\chapters@set{20261126T031100}{41}% theorem-c-and-the-main-inequality
\providecommand\chapnumber[1]{%
  \@ifundefined{chapters@#1}%
    {\PackageError{chapters}{Sheet #1 is not listed in the index}%
      {List it in 20261125T012000--grh-workbook__guide.org\MessageBreak
       and run scripts/make-chapters.py.}}%
    {\edef\thechap{\@nameuse{chapters@#1}}}}
\makeatother
\chapnumber{<sheet ID>}
\makeatletter
\providecommand\chapters@set[2]{\@namedef{chapters@#1}{#2}}
\chapters@set{20261125T001000}{1}% zeros-to-ramp-sums
\chapters@set{20261125T001100}{2}% inert-primes-and-estimate-b
\chapters@set{20261125T005000}{3}% galois-fourier-and-averaging
\chapters@set{20261125T005500}{4}% columns-from-a-sphere
\chapters@set{20261125T000000}{5}% slope-method
\chapters@set{20261126T030100}{6}% a-twist-for-each-prime
\chapters@set{20261125T003000}{7}% the-chain-and-the-square-root
\chapters@set{20261125T010000}{8}% rank-lemma-and-zero-estimate
\chapters@set{20261125T010500}{9}% poincare-series-and-near-orthogonality
\chapters@set{20261124T223000}{10}% elliptic-curves-for-the-proof
\chapters@set{20261125T009500}{11}% fock-model-and-theta-intertwiners
\chapters@set{20261124T223100}{12}% complex-multiplication-and-inert-reduction
\chapters@set{20261125T008000}{13}% one-cm-determinant
\chapters@set{20261124T214500}{14}% supersingular-theta-divisibility
\chapters@set{20261124T214600}{15}% the-isotropic-plane
\chapters@set{20261125T013700}{16}% pairings-on-a-formal-group-of-height-two
\chapters@set{20261125T013500}{17}% hecke-nondegeneracy-on-the-tree
\chapters@set{20261126T030200}{18}% the-bound-is-attained
\chapters@set{20261125T009000}{19}% kronecker-limit-and-chowla-selberg
\chapters@set{20261125T002000}{20}% definite-theta-lifts
\chapters@set{20261125T015500}{21}% siegel-weil-in-the-smallest-case
\chapters@set{20261125T002100}{22}% definite-theta-lifts-and-norms
\chapters@set{20261125T016500}{23}% hecke-operators-and-theta-lifts
\chapters@set{20261125T016700}{24}% divisibility-along-the-chain
\chapters@set{20261125T004000}{25}% denominators-and-sha
\chapters@set{20261125T004100}{26}% tate-twists-and-sha
\chapters@set{20261126T020430}{27}% the-minus-part-by-gauss-sums
\chapters@set{20261125T014100}{28}% congruence-modules
\chapters@set{20261125T014500}{29}% eisenstein-congruences-and-selmer-classes
\chapters@set{20261125T014700}{30}% pseudocharacters-and-the-length-bound
\chapters@set{20261125T011000}{31}% galois-representations-black-box
\chapters@set{20261126T030300}{32}% the-data
\chapters@set{20261126T020210}{33}% bessel-poincare-series-and-their-gram-operators
\chapters@set{20261126T030400}{34}% the-seesaw-identity
\chapters@set{20261126T030500}{35}% divisibility-of-the-composite-lift
\chapters@set{20261126T030600}{36}% theorem-a-short-forms
\chapters@set{20261126T030700}{37}% the-determinant-of-a-chain
\chapters@set{20261126T030800}{38}% the-cost-of-one-lift
\chapters@set{20261126T030900}{39}% denominators-of-projections-to-lifts
\chapters@set{20261126T031000}{40}% twists-with-few-congruences
\chapters@set{20261126T031100}{41}% theorem-c-and-the-main-inequality
\providecommand\chapnumber[1]{%
  \@ifundefined{chapters@#1}%
    {\PackageError{chapters}{Sheet #1 is not listed in the index}%
      {List it in 20261125T012000--grh-workbook__guide.org\MessageBreak
       and run scripts/make-chapters.py.}}%
    {\edef\thechap{\@nameuse{chapters@#1}}}}
\makeatother
\chapnumber{<sheet ID>}
\makeatletter
\providecommand\chapters@set[2]{\@namedef{chapters@#1}{#2}}
\chapters@set{20261125T001000}{1}% zeros-to-ramp-sums
\chapters@set{20261125T001100}{2}% inert-primes-and-estimate-b
\chapters@set{20261125T005000}{3}% galois-fourier-and-averaging
\chapters@set{20261125T005500}{4}% columns-from-a-sphere
\chapters@set{20261125T000000}{5}% slope-method
\chapters@set{20261126T030100}{6}% a-twist-for-each-prime
\chapters@set{20261125T003000}{7}% the-chain-and-the-square-root
\chapters@set{20261125T010000}{8}% rank-lemma-and-zero-estimate
\chapters@set{20261125T010500}{9}% poincare-series-and-near-orthogonality
\chapters@set{20261124T223000}{10}% elliptic-curves-for-the-proof
\chapters@set{20261125T009500}{11}% fock-model-and-theta-intertwiners
\chapters@set{20261124T223100}{12}% complex-multiplication-and-inert-reduction
\chapters@set{20261125T008000}{13}% one-cm-determinant
\chapters@set{20261124T214500}{14}% supersingular-theta-divisibility
\chapters@set{20261124T214600}{15}% the-isotropic-plane
\chapters@set{20261125T013700}{16}% pairings-on-a-formal-group-of-height-two
\chapters@set{20261125T013500}{17}% hecke-nondegeneracy-on-the-tree
\chapters@set{20261126T030200}{18}% the-bound-is-attained
\chapters@set{20261125T009000}{19}% kronecker-limit-and-chowla-selberg
\chapters@set{20261125T002000}{20}% definite-theta-lifts
\chapters@set{20261125T015500}{21}% siegel-weil-in-the-smallest-case
\chapters@set{20261125T002100}{22}% definite-theta-lifts-and-norms
\chapters@set{20261125T016500}{23}% hecke-operators-and-theta-lifts
\chapters@set{20261125T016700}{24}% divisibility-along-the-chain
\chapters@set{20261125T004000}{25}% denominators-and-sha
\chapters@set{20261125T004100}{26}% tate-twists-and-sha
\chapters@set{20261126T020430}{27}% the-minus-part-by-gauss-sums
\chapters@set{20261125T014100}{28}% congruence-modules
\chapters@set{20261125T014500}{29}% eisenstein-congruences-and-selmer-classes
\chapters@set{20261125T014700}{30}% pseudocharacters-and-the-length-bound
\chapters@set{20261125T011000}{31}% galois-representations-black-box
\chapters@set{20261126T030300}{32}% the-data
\chapters@set{20261126T020210}{33}% bessel-poincare-series-and-their-gram-operators
\chapters@set{20261126T030400}{34}% the-seesaw-identity
\chapters@set{20261126T030500}{35}% divisibility-of-the-composite-lift
\chapters@set{20261126T030600}{36}% theorem-a-short-forms
\chapters@set{20261126T030700}{37}% the-determinant-of-a-chain
\chapters@set{20261126T030800}{38}% the-cost-of-one-lift
\chapters@set{20261126T030900}{39}% denominators-of-projections-to-lifts
\chapters@set{20261126T031000}{40}% twists-with-few-congruences
\chapters@set{20261126T031100}{41}% theorem-c-and-the-main-inequality
\providecommand\chapnumber[1]{%
  \@ifundefined{chapters@#1}%
    {\PackageError{chapters}{Sheet #1 is not listed in the index}%
      {List it in 20261125T012000--grh-workbook__guide.org\MessageBreak
       and run scripts/make-chapters.py.}}%
    {\edef\thechap{\@nameuse{chapters@#1}}}}
\makeatother
\chapnumber{20261125T010500}
\renewcommand{\thesection}{\thechap.\arabic{section}}

\theoremstyle{definition}
\newtheorem{problem}{Problem}[section]
\theoremstyle{plain}
\newtheorem{theorem}{Theorem}[section]

\newcommand{\Z}{\mathbb{Z}}
\newcommand{\Q}{\mathbb{Q}}
\newcommand{\R}{\mathbb{R}}
\newcommand{\C}{\mathbb{C}}
\newcommand{\HH}{\mathbb{H}}
\newcommand{\ii}{\mathrm{i}}
\newcommand{\Ac}{\mathcal{A}}
\newcommand{\Hc}{\mathcal{H}}
\newcommand{\Pt}{\widetilde{P}}
\newcommand{\vol}{\operatorname{vol}}

\setlength{\emergencystretch}{2em}

\title{Poincar\'e series and near orthogonality}
\author{}
\date{November 25, 2026}

\begin{document}
\maketitle
\refstepcounter{chap}\label{chap}

\paragraph{Prerequisites.}
Modular forms for \(\mathrm{SL}_2(\Z)\): \(q\)-expansions, the cusp form \(\Delta\) and the dimension formula \cite[Ch.~VII, \S3]{SeA}.  The residue theorem and the mean value property of holomorphic functions.  The power series \(J_\nu(x) = \sum_{j \ge 0} (-1)^j (x/2)^{\nu + 2j}/(j!\,\Gamma(\nu + j + 1))\).  Gershgorin's theorem on the eigenvalues of a matrix.  Chapter~\ref{rl-chap}, \emph{Few translates span}, for the closing paragraphs only.

Chapter~\ref{rl-chap} asked for a map into a space of automorphic forms that is almost an isometry.  If the weight is large, Poincar\'e series of small index are almost orthogonal, and Petersson's formula proves it with nothing more than a bound for a Bessel function.  This chapter works it out on \(\mathrm{SL}_2(\Z)\), where the vectors are holomorphic Poincar\'e series and their Petersson inner products recover Fourier coefficients up to an explicit factor.

Throughout, \(k \ge 4\) is even, \(\Gamma := \mathrm{SL}_2(\Z)\), \(\Gamma_\infty := \{\pm\begin{pmatrix}1 & n\\ 0 & 1\end{pmatrix} \mid n\in\Z\}\), \(j(\gamma,z) := cz + d\) for \(\gamma = \begin{pmatrix}a & b\\ c & d\end{pmatrix}\), \(e(z) := e^{2\pi\ii z}\), and \(d\nu := y^{-2}\,dx\,dy\) on the upper half-plane \(\HH\).

\section{Poincar\'e series}\label{sec:poincare}

Let \(S_k\) be the space of cusp forms of weight \(k\) for \(\Gamma\), with the Petersson inner product \(\langle f, g\rangle := \int_{\Gamma\backslash\HH} f(z)\overline{g(z)}\,y^k\,d\nu(z)\).  For \(f \in S_k\) write \(f(z) = \sum_{n \ge 1} a_f(n)\,e(nz)\).  For each positive integer \(m\) the \emph{Poincar\'e series}
\[
  P_m(z) := \sum_{\gamma \in \Gamma_\infty\backslash\Gamma} j(\gamma,z)^{-k}\, e(m\gamma z)
\]
converges absolutely and uniformly on compact sets, and \(P_m \in S_k\) \cite[\S\S3.1--3.2]{Iw}.

\begin{problem}\label{pr:unfold}
For \(f \in S_k\) and \(m \ge 1\),
\[
  \langle f, P_m\rangle = \frac{\Gamma(k-1)}{(4\pi m)^{k-1}}\, a_f(m) .
\]
\end{problem}

\begin{problem}\label{pr:span}
The series \(P_m\), \(m \ge 1\), span \(S_k\).
\end{problem}

So \(P_m = 0\) for \(k \in \{4, 6, 8, 10, 14\}\), where \(S_k = 0\), although no term of the series vanishes.  The first nonzero case is \(k = 12\), where \(S_{12} = \C\Delta\) with \(\Delta = \sum_{n \ge 1}\tau(n)\,e(nz)\).

\begin{problem}\label{pr:k12}
For \(k = 12\),
\[
  P_m = \frac{10!\;\tau(m)}{(4\pi m)^{11}\,\lVert\Delta\rVert^2}\,\Delta .
\]
In particular \(P_1\) and \(P_2\) are proportional.
\end{problem}

So distinct Poincar\'e series need not be independent.  Petersson's formula, in the next section, measures how far they are from orthogonal.

\section{Petersson's formula}\label{sec:petersson}

What is the Fourier expansion of \(P_m\)?  For integers \(m, n\) and \(c \ge 1\) let
\[
  S(m,n;c) := \sum_{\substack{d \bmod c\\ (d,c) = 1}} e\Bigl(\frac{m\bar d + nd}{c}\Bigr), \qquad d\bar d \equiv 1 \pmod c,
\]
be the Kloosterman sum.  It is real, and \(S(m,n;c) = S(n,m;c)\).

\begin{problem}\label{pr:cosets}
The map \(\gamma \mapsto (c,d)\) induces a bijection from \(\Gamma_\infty\backslash\Gamma\) onto the set of pairs of coprime integers modulo \(\pm1\).  If \(c \ne 0\), then \(\gamma z = \frac ac - \frac{1}{c(cz+d)}\).
\end{problem}

\begin{problem}\label{pr:lipschitz}
Let \(\ell \ge 2\) be an integer, \(\beta \in \R\) and \(y > 0\).  Then
\[
  \int_{\ii y - \infty}^{\ii y + \infty} z^{-\ell}\, e(-\beta z)\, dz =
  \begin{cases} (-2\pi\ii)^\ell\, \beta^{\ell-1}/(\ell-1)! & \beta > 0,\\ 0 & \beta \le 0. \end{cases}
\]
\end{problem}

\begin{problem}\label{pr:bessel-int}
Let \(\alpha > 0\), \(\beta \in \R\) and \(y > 0\).  Then \(\int_{\ii y - \infty}^{\ii y + \infty} z^{-k}\, e(-\alpha/z - \beta z)\, dz\) vanishes for \(\beta \le 0\), and for \(\beta > 0\) equals
\[
  2\pi\ii^{-k} \Bigl(\frac\beta\alpha\Bigr)^{(k-1)/2} J_{k-1}\bigl(4\pi\sqrt{\alpha\beta}\bigr) .
\]
\end{problem}

\begin{problem}\label{pr:fourier}
\(P_m(z) = \sum_{n \ge 1} p_m(n)\, e(nz)\), where
\[
  p_m(n) = \delta(m,n) + 2\pi\ii^{-k} \Bigl(\frac nm\Bigr)^{(k-1)/2} \sum_{c \ge 1} \frac{S(m,n;c)}{c}\, J_{k-1}\Bigl(\frac{4\pi\sqrt{mn}}{c}\Bigr) .
\]
\end{problem}

Problems~\ref{pr:unfold} and~\ref{pr:fourier} together give Petersson's formula \cite{Pe}, \cite[\S3.2]{Iw}, \cite[\S2]{ILS}.  Let
\[
  T_k(m,n) := 2\pi\ii^{-k} \sum_{c \ge 1} \frac{S(m,n;c)}{c}\, J_{k-1}\Bigl(\frac{4\pi\sqrt{mn}}{c}\Bigr) ,
\]
a real symmetric function of \(m, n \ge 1\).

\begin{theorem}[Petersson]\label{thm:petersson}
For every orthonormal basis \((f_j)\) of \(S_k\) and all \(m, n \ge 1\),
\[
  \frac{\Gamma(k-1)}{(4\pi\sqrt{mn})^{k-1}} \sum_j \overline{a_{f_j}(m)}\, a_{f_j}(n) = \delta(m,n) + T_k(m,n) .
\]
\end{theorem}

\begin{problem}\label{pr:vanishing}
For \(k \in \{4, 6, 8, 10, 14\}\) and all \(m, n \ge 1\),
\[
  \sum_{c \ge 1} \frac{S(m,n;c)}{c}\, J_{k-1}\Bigl(\frac{4\pi\sqrt{mn}}{c}\Bigr) = -\frac{\ii^k}{2\pi}\,\delta(m,n) .
\]
\end{problem}

Numerically, for \(m = n = 1\) the partial sums over \(c \le 800\) reproduce these values to within \(3 \cdot 10^{-8}\) for \(k = 4\) and \(10^{-13}\) for \(k \ge 6\).  For \(k = 12\) one finds \(1 + T_{12}(1,1) = 2.84028\ldots\), which by Theorem~\ref{thm:petersson} is \(10!/\bigl((4\pi)^{11}\lVert\Delta\rVert^2\bigr)\), so that \(\lVert\Delta\rVert^2 = 1.03536\ldots \cdot 10^{-6}\).

\section{Near orthogonality}\label{sec:near}

Normalize the Poincar\'e series by
\[
  \Pt_m := \frac{(4\pi m)^{(k-1)/2}}{\Gamma(k-1)^{1/2}}\, P_m ,
\]
and for each positive integer \(M\) let \(G_M\) be the \(M \times M\) Gram matrix \(\bigl(\langle \Pt_m, \Pt_n\rangle\bigr)_{1 \le m,n \le M}\).  We use the operator norm on \(\C^M\) with its Euclidean inner product.

\begin{problem}\label{pr:gram}
\(\langle \Pt_m, \Pt_n\rangle = \delta(m,n) + T_k(m,n)\).
\end{problem}

So \(G_M\) is the identity plus the matrix of Kloosterman terms, and the question is how small \(J_{k-1}\) is.  We use Poisson's integral
\[
  J_\nu(x) = \frac{(x/2)^\nu}{\sqrt\pi\,\Gamma(\nu + \frac12)} \int_{-1}^1 (1 - t^2)^{\nu - 1/2} \cos(xt)\, dt \qquad (\nu > -\tfrac12,\ x > 0),
\]
see \cite[(10.9.4)]{DLMF}.

\begin{problem}\label{pr:J-bound}
For \(\nu > -\frac12\) and \(x > 0\), \(\lvert J_\nu(x)\rvert \le (x/2)^\nu/\Gamma(\nu + 1)\).
\end{problem}

\begin{problem}\label{pr:T-bound}
\(\lvert T_k(m,n)\rvert \le 2\pi\,\zeta(k-1)\,(2\pi\sqrt{mn})^{k-1}/(k-1)!\).
\end{problem}

\begin{problem}\label{pr:gershgorin}
If \(1 \le M \le (k-1)/(4\pi e)\), then every eigenvalue of \(G_M\) lies within \(k\,2^{-k}\) of \(1\).
\end{problem}

Problems~\ref{pr:gram}--\ref{pr:gershgorin} prove the following.

\begin{theorem}\label{thm:near}
If \(1 \le M \le (k-1)/(4\pi e)\), then \(\lVert G_M - 1\rVert \le k\,2^{-k}\).  In particular \(P_1, \dots, P_M\) are linearly independent, and \(\dim S_k \ge \lfloor (k-1)/(4\pi e)\rfloor\).
\end{theorem}

Here \(e = 2.718\ldots\) is Euler's number, and \(4\pi e = 34.15\ldots\).  The threshold is not the true one.  For \(0<x\le r\nu\) with \(0<r<1\) fixed, the function \(J_\nu(x)\) is exponentially small in \(\nu\), and Siegel's inequality makes this precise: for \(\nu \ge 0\) and \(0 < x \le 1\),
\[
  \lvert J_\nu(\nu x)\rvert \le \Bigl(\frac{x\, e^{\sqrt{1 - x^2}}}{1 + \sqrt{1 - x^2}}\Bigr)^{\nu}
\]
\cite[(10.14.5)]{DLMF}.

\begin{problem}\label{pr:siegel}
Let \(0 < \theta < 1/(4\pi)\).  There are constants \(0 < \varrho < 1\) and \(C > 0\), depending only on \(\theta\), such that \(\lVert G_M - 1\rVert \le C k\varrho^k\) for every even \(k \ge 4\) and every integer \(M\) with \(1 \le M \le \theta(k-1)\).
\end{problem}

So for large \(k\) the \(\Pt_m\) with \(m \le \theta(k-1)\), for a fixed \(\theta < 1/(4\pi)\), are almost orthonormal.  For \(m\) near \((k-1)/(4\pi)\) the argument \(4\pi m\) of the term \(c = 1\) of \(T_k(m,m)\) reaches the transition region \(x \approx \nu\) of \(J_\nu\), where \(J_\nu(x)\) is no longer exponentially small but of size up to about \(\nu^{-1/3}\), and the diagonal entry is in general no longer exponentially close to \(1\).  For \(k = 100\), where \((k-1)/(4\pi) = 7.87\ldots\) and \(\dim S_{100} = 8\), a computation of the Kloosterman terms gives the following.  The eigenvalues of \(G_7\) lie in \([0.99999, 1.0094]\).  Those of \(G_8\) lie in \([0.982, 1.808]\), so \(P_1,\dots,P_8\) form a basis of \(S_{100}\), and the Gram matrix of the normalized basis \(\Pt_1,\dots,\Pt_8\) has an eigenvalue near \(1.808\).  \(G_9\) is singular, as it must be.

\section{The dimension as formal degree times covolume}\label{sec:dim}

How large a proportion of \(S_k\) does Theorem~\ref{thm:near} fill?  Let \(\Gamma\backslash\HH\) carry the measure \(\nu\).

\begin{problem}\label{pr:vol}
\(\nu(\Gamma\backslash\HH) = \pi/3\).
\end{problem}

\begin{problem}\label{pr:dim}
\(\bigl\lvert \dim S_k - \frac{k-1}{12}\bigr\rvert \le \frac{13}{12}\) for every even \(k \ge 4\).
\end{problem}

For fixed \(0<\theta<1/(4\pi)\), the first \(\lfloor\theta(k-1)\rfloor\) series in Problem~\ref{pr:siegel} span a proportion of \(S_k\) tending to \(12\theta\) as \(k\) tends to infinity through even weights.  These limits approach \(3/\pi=0.954\ldots\) as \(\theta\) increases to \(1/(4\pi)\).  The full initial segment in Theorem~\ref{thm:near} gives limiting proportion \(3/(\pi e)=0.351\ldots\).  Where does \((k-1)/12\) come from?  It is the product of \(\pi/3\) and a constant attached to the weight alone.

Let \(G:=\mathrm{PSL}_2(\R)\) act on \(\HH\), and normalize Haar measure \(dg\) by \(\int_G F(g\ii)\,dg=\int_\HH F\,d\nu\) for compactly supported continuous \(F\).  With Haar probability on the stabilizer \(\mathrm{PSO}(2)\) in the quotient integral formula, \(\vol(\mathrm{PSL}_2(\Z)\backslash G)=\pi/3\).  Let \(\Ac_k\) be the space of holomorphic functions \(f\) on \(\HH\) with \(\lVert f\rVert^2 := \int_\HH \lvert f\rvert^2 y^k\,d\nu < \infty\), on which \(G\) acts unitarily by \((\pi(g)f)(z) := f(g^{-1}z)\,j(g^{-1},z)^{-k}\).  This is the holomorphic discrete series of weight \(k\).  For \(w \in \HH\) put
\[
  h_w(z) := \Bigl(\frac{z - \bar w}{2\ii}\Bigr)^{-k} .
\]
For a unit vector \(u\) of an irreducible unitary representation \(\pi\) with square-integrable matrix coefficients, the \emph{formal degree} \(d_\pi\) is defined by \(\int_G \lvert\langle\pi(g)u, u\rangle\rvert^2\,dg = 1/d_\pi\).  It does not depend on \(u\), by the Schur orthogonality relations for such representations \cite{HC}.  We compute it with \(u\) proportional to \(h_\ii\).

\begin{problem}\label{pr:cayley}
\(h_w \in \Ac_k\), and \(\langle f, h_\ii\rangle = \frac{4\pi}{k-1}\, f(\ii)\) for every \(f \in \Ac_k\).
\end{problem}

\begin{problem}\label{pr:covariance}
\(\pi(g)h_w = \overline{j(g,w)}^{\,-k}\, h_{gw}\) for \(g \in \mathrm{SL}_2(\R)\) and \(w \in \HH\).
\end{problem}

\begin{problem}\label{pr:formal-degree}
Let \(u := h_\ii/\lVert h_\ii\rVert\), and \(w := g\ii\).  Then \(\lvert\langle\pi(g)u, u\rangle\rvert^2 = \bigl(4\operatorname{Im} w/\lvert w + \ii\rvert^2\bigr)^k\), and the formal degree of \(\pi\) is \((k-1)/(4\pi)\).
\end{problem}

With Problem~\ref{pr:vol} this gives
\[
  d_\pi \cdot \vol(\mathrm{PSL}_2(\Z)\backslash G) = \frac{k-1}{4\pi}\cdot\frac\pi3 = \frac{k-1}{12} ,
\]
which is \(\dim S_k\) up to the bounded error of Problem~\ref{pr:dim}.  This is not a coincidence.  As in Problem~\ref{pr:unfold}, and using Problems~\ref{pr:cayley} and~\ref{pr:covariance}, one checks that the reproducing kernel \(B(z,w) := \sum_j f_j(z)\overline{f_j(w)}\) of \(S_k\), for an orthonormal basis \((f_j)\), is
\[
  B(z,w) = \frac{k-1}{4\pi} \sum_{\gamma \in \mathrm{PSL}_2(\Z)} j(\gamma,z)^{-k}\, h_w(\gamma z) .
\]
Since \(\dim S_k = \int_{\Gamma\backslash\HH} B(z,z)\,y^k\,d\nu(z)\), the term \(\gamma = 1\) contributes exactly the formal degree times the covolume, and by Problem~\ref{pr:dim} all other terms together contribute a bounded amount.  For a compact unitary group with Haar probability, the formal degree of an irreducible representation is its dimension.  In the definite example of Chapter~\ref{dt-chap}, for rank \(n\ge1\) and weight \(\lambda\) with \(4\mid\lvert\lambda\rvert\), the central scalars contribute \(4M_n\dim V_\lambda\), where \(M_n\) is the mass (Problem~\ref{dt-pr:dim-forms}).

Theorem~\ref{thm:near} says that the map \(\C^M \to S_k\), \(e_m \mapsto \Pt_m\), is almost an isometry when \(M\) is at most a fixed proportion of \(\dim S_k\).  Its image therefore occupies only part of \(S_k\).  Problem~\ref{pr:siegel} extends near orthonormality to a longer initial segment, and Problem~\ref{pr:span} shows that the full family spans \(S_k\).  In the setting of Chapter~\ref{rl-chap}, a finite group acts on the ambient space.  A near-isometry estimate for the projected map and a second-moment bound for the Gram operators of the unprojected map and its translates give a lower bound for the dimension spanned by a few projected translates in the target (Theorems~\ref{rl-thm:spread} and~\ref{rl-thm:few}).

\paragraph{The analogue in the rewritten Lean proof.}
This chapter gives a classical model for the estimates in Chapter~\ref{bp-chap}.  In the rewritten Lean proof \cite{GL}, the principal-level forms on the unitary group of a hermitian space \(Y_0\) over a CM field \(E\) play the role of \(S_k\).  Bessel periods play the role of Fourier coefficient extraction and are represented by a Poincar\'e map.  Its values on the source coordinate basis are Bessel--Poincar\'e series.  The two estimates below concern the map followed by the projector onto the good forms and the dimension of those forms.  For the definition of Bessel periods see \cite[\S\S12, 23]{GGP}.

\paragraph{In the Lean formalization.}
Fix admissible data \(D\) of the rewritten Lean proof \cite{GL}.  Write \(E\) for \texttt{Data.E}, \(Y_0\) for \texttt{D.spaceY 0}, and \(r\) for the rank of \(Y_0\) over \(E\).  Write \(S\) for the source forms \texttt{Data.SourceForm}, with coordinates \texttt{Data.sourceFormCoords}, and \(D_X\) for their complex dimension.  The Poincar\'e map \texttt{Data.poincare} represents the Bessel periods at the chosen point and linear form by \texttt{Data.inner\_poincareAt}.  By \texttt{Data.poincare\_eq\_poincareMap\_comp} and \texttt{Filling.Fits.coe\_poincareMap}, it takes the coordinate basis to the corresponding Bessel--Poincar\'e series.

Let \(P\) be \texttt{Data.scaledPoincare}, the Poincar\'e map in these Euclidean coordinates multiplied by the positive factor \texttt{Data.poincareScale}.  Its codomain is the space \texttt{Data.principalForms} on \(Y_0\), of the specified discrete-series type at infinity and principal level at the finite places.  Let \(e_0\) denote the action of \texttt{D.e\(_0\)} on this space.  It projects onto the space \(T\) of good forms (\texttt{Data.targetForms}), of complex dimension \(d_T\).  Let \(b\) be the weight \texttt{D.b} at the designated infinite places.  For every complex subspace \(U\) of the source coordinate space, \texttt{Data.isNearIsometry\_e\(_0\)\_comp\_scaledPoincare} gives
\[
  \lVert e_0 P\,\mathrm{pr}_U\rVert_{\mathrm{HS}}^2
  \ge \dim U - 5b^{-4}D_X ,
\]
where \(\mathrm{pr}_U\) is the orthogonal projection and \(\lVert\cdot\rVert_{\mathrm{HS}}\) is the Hilbert--Schmidt norm.

By \texttt{Data.targetDim\_le\_card\_mul\_besselDegree\_mul\_mass},
\[
  d_T \le \#\bar G \cdot \deg \cdot \mathrm{mass}(E,r) .
\]
Here \(\bar G\) is \texttt{Data.chosenGL}, the product of the residue general linear groups at the chosen places above the two level primes, \(\deg\) is the product over all infinite places of the degrees of the local Bessel systems, and \(\mathrm{mass}(E,r)\) is the mass of the genus of the standard hermitian lattice \(\mathcal O_E^r\).  Both inequalities are in \texttt{grh/src/New/Spanning/SecondMoment/SpanningTheorem.lean}.

\section*{Hints}

\subsection*{\S\ref{sec:poincare}}
\begin{description}
\item[\ref{pr:unfold}] Unfold the sum over \(\Gamma_\infty\backslash\Gamma\) to the strip \(0 \le x \le 1\), \(y > 0\), and integrate in \(x\) first.
\item[\ref{pr:span}] Problem~\ref{pr:unfold}.
\item[\ref{pr:k12}] Expand \(P_m\) in the basis \(\Delta\) and use Problem~\ref{pr:unfold}.  Note that \(\tau(1) = 1\), \(\tau(2) = -24\).
\end{description}

\subsection*{\S\ref{sec:petersson}}
\begin{description}
\item[\ref{pr:cosets}] The bottom row of a matrix in \(\Gamma\) is coprime, and two matrices with the same bottom row differ by left multiplication by an element of \(\Gamma_\infty\).
\item[\ref{pr:lipschitz}] For \(\beta\ne0\), close the contour in the half-plane where \(e(-\beta z)\) decays.  For \(\beta=0\), use a primitive of \(z^{-\ell}\).
\item[\ref{pr:bessel-int}] Expand \(e(-\alpha/z)\) in its power series and apply Problem~\ref{pr:lipschitz} term by term.  For \(\beta \le 0\), \(\lvert e(-\alpha/z)\rvert \le 1\) in \(\HH\).
\item[\ref{pr:fourier}] The terms with \(c > 0\) and fixed \(c\) are indexed by \(d \in \Z\) coprime to \(c\); group them by \(d \bmod c\).  Each group is a sum over translates \(z \mapsto z + r\), \(r \in \Z\), so its \(n\)-th Fourier coefficient is an integral over a whole horizontal line.  Use Problems~\ref{pr:cosets} and~\ref{pr:bessel-int}.
\item[\ref{pr:vanishing}] Theorem~\ref{thm:petersson}.
\end{description}

\subsection*{\S\ref{sec:near}}
\begin{description}
\item[\ref{pr:gram}] Problem~\ref{pr:unfold} with \(f = P_m\), and Problem~\ref{pr:fourier}.
\item[\ref{pr:J-bound}] Bound \(\lvert\cos(xt)\rvert\) by \(1\) and evaluate the Beta integral.
\item[\ref{pr:T-bound}] \(\lvert S(m,n;c)\rvert \le c\).
\item[\ref{pr:gershgorin}] \((k-1)! \ge ((k-1)/e)^{k-1}\).  Then Gershgorin's theorem.
\item[\ref{pr:siegel}] Treat the term \(c = 1\) with Siegel's inequality.  The base on its right is increasing in \(x\) and equals \(1\) at \(x = 1\).  For the terms \(c \ge 2\), use Problem~\ref{pr:J-bound}.
\end{description}

\subsection*{\S\ref{sec:dim}}
\begin{description}
\item[\ref{pr:vol}] The standard fundamental domain \(\lvert x\rvert \le \frac12\), \(\lvert z\rvert \ge 1\); integrate in \(y\) first.
\item[\ref{pr:dim}] Write \(k = 12j + r\) with \(0 \le r < 12\) and use the dimension formula.
\item[\ref{pr:cayley}] Put \(\zeta := (z - \ii)/(z + \ii)\) and \(F(\zeta) := f(z)\bigl((z+\ii)/2\ii\bigr)^k\).  Express \(\lvert f\rvert^2 y^k\,d\nu\) in terms of \(F\) and \(\zeta\), and use the mean value property on the circles \(\lvert\zeta\rvert = r\).  For \(h_w \in \Ac_k\), use \(4\operatorname{Im}z\operatorname{Im}w/\lvert z - \bar w\rvert^2 = 1 - \lvert\zeta_w\rvert^2\) with \(\zeta_w := (z - w)/(z - \bar w)\).
\item[\ref{pr:covariance}] \(gz - \overline{gw} = (z - \bar w)/\bigl(j(g,z)\overline{j(g,w)}\bigr)\).
\item[\ref{pr:formal-degree}] Problems~\ref{pr:cayley} and~\ref{pr:covariance}.  For the integral, use geodesic polar coordinates about \(\ii\), in which \(4\operatorname{Im} w/\lvert w + \ii\rvert^2 = \cosh^{-2}(\rho/2)\).
\end{description}

\begin{thebibliography}{GGP}
\bibitem[Alp]{Alp} L.~Alp\"oge, untitled note announcing a proof of the Riemann hypothesis, 20 November 2026, \url{https://alpo.ge/remark.pdf}.
\bibitem[DLMF]{DLMF} F.~W.~J.~Olver, D.~W.~Lozier, R.~F.~Boisvert and C.~W.~Clark (eds.), \emph{NIST Handbook of Mathematical Functions}, Cambridge Univ.\ Press, 2010; \url{https://dlmf.nist.gov}, Ch.~10.
\bibitem[GGP]{GGP} W.~T.~Gan, B.~H.~Gross and D.~Prasad, \emph{Symplectic local root numbers, central critical \(L\)-values, and restriction problems in the representation theory of classical groups}, Ast\'erisque \textbf{346} (2012), 1--109.
\bibitem[GL]{GL} L.~Sernau, \emph{GRH for Dirichlet \(L\)-functions, in Lean~4}, a shorter rewriting of the Lean proof, text by Claude, \url{https://github.com/lsernau/grh-lean}, commit \texttt{72fead1}, readable edition \texttt{grh/}.
\bibitem[HC]{HC} Harish-Chandra, \emph{Harmonic analysis on semisimple Lie groups}, Bull.\ Amer.\ Math.\ Soc.\ \textbf{76} (1970), 529--551.
\bibitem[ILS]{ILS} H.~Iwaniec, W.~Luo and P.~Sarnak, \emph{Low lying zeros of families of \(L\)-functions}, Publ.\ Math.\ IH\'ES \textbf{91} (2000), 55--131.
\bibitem[Iw]{Iw} H.~Iwaniec, \emph{Topics in Classical Automorphic Forms}, Grad.\ Stud.\ Math.\ \textbf{17}, Amer.\ Math.\ Soc., Providence, 1997.
\bibitem[Pe]{Pe} H.~Petersson, \emph{\"Uber die Entwicklungskoeffizienten der automorphen Formen}, Acta Math.\ \textbf{58} (1932), 169--215.
\bibitem[SeA]{SeA} J.-P.~Serre, \emph{A Course in Arithmetic}, Grad.\ Texts in Math.\ 7, Springer, 1973.
\end{thebibliography}

\end{document}
